% --- Chapter: Object-Oriented Programming ---
\section{Object-Oriented Programming}

CRISP v0.1.5 introduces Python/Perl-style object-oriented programming with classes, constructors, inheritance, and static methods. OOP in CRISP is prototype-free and class-based, with clean syntax modeled on modern dynamic languages.

\subsection{Defining a Class}

A class is defined with the \texttt{class} keyword. Methods are defined with \texttt{fn} inside the class body. The special \texttt{new} method is the constructor:

\begin{tcolorbox}[colback=codebg, colframe=darkpurple, colbacktitle=darkpurple, title=\textbf{\textcolor{codegold}{Class Definition}}, fonttitle=\bfseries, sharp corners]
\begin{lstlisting}
class Point {
    fn new(x, y) {
        self.x = x;
        self.y = y;
    }

    fn distance() {
        return sqrt(pow(self.x, 2) + pow(self.y, 2));
    }

    fn describe() {
        return "Point(" + self.x + ", " + self.y + ")";
    }
}
\end{lstlisting}
\end{tcolorbox}

\subsection{Creating Objects}

Objects are created with the \texttt{new} keyword, which calls the constructor:

\begin{tcolorbox}[colback=codebg, colframe=darkpurple, colbacktitle=darkpurple, title=\textbf{\textcolor{codegold}{Creating Objects}}, fonttitle=\bfseries, sharp corners]
\begin{lstlisting}
let p = new Point(3, 4);

say p.x;           # 3
say p.y;           # 4
say p.distance();  # 5
say p.describe();  # Point(3, 4)

# Mutate fields directly
p.x = 10;
say p.distance();  # sqrt(100 + 16) = 10.77...
\end{lstlisting}
\end{tcolorbox}

\subsection{Class.method() Auto-Instantiation}

Calling a method on a class (not an instance) auto-creates an object. This is syntactic sugar that calls \texttt{new} with the arguments, then calls the method:

\begin{tcolorbox}[colback=codebg, colframe=darkpurple, colbacktitle=darkpurple, title=\textbf{\textcolor{codegold}{Auto-Instantiation}}, fonttitle=\bfseries, sharp corners]
\begin{lstlisting}
# These are equivalent:
let p1 = Point.new(5, 12);  # explicit new
let p2 = new Point(5, 12);  # contextual new keyword

say p1.distance();  # 13
say p2.distance();  # 13

# Auto-instantiation works for ANY method, not just new:
let p3 = Point.move_to(6, 8);  # creates Point(0,0), moves to (6,8)
say p3.describe();  # Point(6, 8)
\end{lstlisting}
\end{tcolorbox}

\subsection{Inheritance}

Classes can extend other classes with the \texttt{extends} keyword. Child classes inherit all methods and fields from the parent:

\begin{tcolorbox}[colback=codebg, colframe=darkpurple, colbacktitle=darkpurple, title=\textbf{\textcolor{codegold}{Inheritance}}, fonttitle=\bfseries, sharp corners]
\begin{lstlisting}
class Point3D extends Point {
    fn new(x, y, z) {
        super(x, y);      # call parent constructor
        self.z = z;
    }

    fn distance() {
        let d2d = super.distance();  # call parent method
        return sqrt(pow(d2d, 2) + pow(self.z, 2));
    }

    fn describe() {
        return "Point3D(" + self.x + ", " + self.y +
               ", " + self.z + ")";
    }
}

let p3 = new Point3D(3, 4, 12);
say p3.describe();  # Point3D(3, 4, 12)
say p3.distance();  # sqrt(3²+4²+12²) = 13
\end{lstlisting}
\end{tcolorbox}

\subsection{Super Calls}

\texttt{super(args...)} calls the parent class constructor. \texttt{super.method(args...)} calls a parent method:

\begin{tcolorbox}[colback=codebg, colframe=darkpurple, colbacktitle=darkpurple, title=\textbf{\textcolor{codegold}{Super Calls}}, fonttitle=\bfseries, sharp corners]
\begin{lstlisting}
class Animal {
    fn new(name, age) {
        self.name = name;
        self.age = age;
    }

    fn speak() {
        say self.name + " makes a sound";
    }

    fn birthday() {
        self.age = self.age + 1;
        say self.name + " turned " + self.age + "!";
    }
}

class Dog extends Animal {
    fn new(name, age, breed) {
        super(name, age);       # Animal.new(name, age)
        self.breed = breed;
    }

    fn speak() {
        say self.name + " barks!";  # override
    }
}

let rex = new Dog("Rex", 3, "German Shepherd");
rex.speak();      # Rex barks!
rex.birthday();   # Rex turned 4!  (inherited from Animal)
\end{lstlisting}
\end{tcolorbox}

\subsection{Static Methods}

Static methods are defined with \texttt{static fn} and are called on the class (not instances). They don't have access to \texttt{self}:

\begin{tcolorbox}[colback=codebg, colframe=darkpurple, colbacktitle=darkpurple, title=\textbf{\textcolor{codegold}{Static Methods}}, fonttitle=\bfseries, sharp corners]
\begin{lstlisting}
class MathUtil {
    static fn square(x) { return x * x; }
    static fn cube(x)   { return x * x * x; }
}

say MathUtil.square(7);  # 49
say MathUtil.cube(3);    # 27

# Static methods can't be called on instances
let m = new MathUtil();
# m.square(5);  # Error — static methods are class-only
\end{lstlisting}
\end{tcolorbox}

\begin{tcolorbox}[colback=lightlavender, colframe=softvioletdark, title=\textbf{\textcolor{darkpurple}{Design Note: OOP Implementation}}, fonttitle=\bfseries, sharp corners]
CRISP's OOP system is implemented as a thin layer over hashes:

\begin{itemize}
    \item \textbf{Classes:} Stored as \texttt{Value::Class \{ name, parent, methods, static\_methods \}} in the environment. A class is a value like any other—it can be assigned to variables, passed to functions, and inspected.
    \item \textbf{Objects:} Stored as \texttt{Value::Object \{ class, fields, methods \}}. The object carries a copy of its class's methods (snapshotted at construction time) and a reference to the class name for inheritance chain lookups.
    \item \textbf{Inheritance:} Method resolution walks the parent chain at call time. \texttt{super} calls skip the current class and start at the parent, preventing infinite recursion.
    \item \textbf{Self:} Bound automatically by \texttt{call\_method} when dispatching. \texttt{self} is a special variable injected into the method's environment.
\end{itemize}

This design keeps the interpreter simple while providing the full OOP experience—constructors, inheritance, method dispatch, and static methods—all built on CRISP's existing hash and function primitives.
\end{tcolorbox}
